This paper has proposed and validated a \textbf{local-geometry-aware contact modeling strategy} for complementarity-based rigid contact. Within an existing hard-contact solution framework, local SDF information is used to construct contact representations near the active region, while \texttt{SlidingPatch} and \texttt{CompactDiscrete} are selected according to the contact state. The method does not replace the underlying time-stepping or LCP solution chain; instead, it strengthens local geometric representation at the contact-modeling level. The regime choice is always preset at the case or application level. The reported results therefore support a modeling principle rather than a runtime automatic regime-identification algorithm.

The numerical results show that local geometry awareness is effective within the intended scope of complementarity-based rigid contact, but that its benefit is strongly state-dependent. For sustained sliding with relatively smooth local geometry, continuous patch representations make fuller use of local geometric information and improve trajectory and velocity responses while keeping the additional modeling cost modest, thereby preserving computational efficiency and improving contact-response accuracy at the same time. For contact onset, state transitions, and locally singular configurations, more compact and conservative discrete representations provide better robustness. The continuous-contact benchmarks therefore validate the full local-geometry-aware modeling stack, whereas the difference between first- and second-order SDF isolates the incremental contribution of curvature feedforward. The industrial cam, analytical roller, and Onset Stress cases support this interpretation, while the Twin Gear benchmark shows that a structural error floor remains in nonsmooth meshing with violent branch switching.

The main conclusion can therefore be stated as follows: \textbf{within the intended application range of complementarity-based rigid contact, a regime-matched local geometric representation can maintain computational efficiency while achieving more accurate contact responses without modifying the underlying solution chain.} Geometry awareness becomes a stable and useful numerical ingredient only when the representation is coordinated with the contact state. The present results do not support a claim of universal optimality for any single regime. They instead support a modeling principle in which local representations are configured according to contact state under clearly stated applicability conditions and failure boundaries. This conclusion provides a clearer basis for the construction, comparison, and engineering deployment of future complementarity-based rigid-contact models.

From a modeling perspective, these results indicate that, even when the underlying LCP solution chain is left unchanged, substantial room for improvement remains at the geometry-input layer of rigid-contact computation. In cases where response error is dominated by gap-prediction bias, distortion of the local normal manifold, and noise during contact activation, the way local geometric information enters the contact model directly affects final accuracy. The present contribution is therefore better understood as a geometric-modeling reinforcement of an established complementarity-based rigid-contact pipeline: its main role is not to replace the underlying solution theory, but to improve the quality of the existing modeling interface through low-intrusion local geometric representations.

Although the present method has been validated only for complementarity-based rigid contact, the current results already indicate that the local geometric and differential information encoded by the SDF can serve as an effective auxiliary quantity in contact computation, improving numerical response under clearly defined applicability conditions and indicating that computational efficiency and response accuracy can coexist within the intended application regimes. In this sense, local geometry awareness should not be viewed solely as a case-specific refinement within the current rigid-body framework, but may point toward a broader modeling direction for contact computation. Future work will therefore examine how richer local spatial geometric descriptors can be incorporated systematically into contact computation for flexible-body contact and for broader classes of contact solvers.
